\subsubsection{Rueda de molino}

Una barra de longitud \( L \) con una rueda de radio \( R \) unida a uno de sus extremos está rotando alrededor del eje vertical \( OA \) con una velocidad angular constante \( \Omega \) relativa a un marco de referencia fijo, como se muestra en la Figura~\ref{fig:rueda_brazo_final}. La rueda está vertical y rueda sin deslizamiento a lo largo de una superficie horizontal fija.\\

Determine la velocidad angular y la aceleración angular de la rueda por un observador en un marco de referencia fijo.\\

\begin{figure}[!ht]
    \centering
    % 1. Ajuste de Perspectiva (Ángulo de elevación y azimut)
    \tdplotsetmaincoords{60}{130}

    \begin{tikzpicture}[tdplot_main_coords, scale=1.5, >=Latex]

        % --- Parámetros Ajustados (Proporciones más realistas) ---
        \def\L{3.5}      % Longitud del brazo
        \def\R{1}      % Radio del disco
        \def\Hshaft{2.0} % Altura del eje central
        \def\zO{1}     % Altura de conexión del brazo
        \def\DiskThick{0.2} % Grosor del disco

        % --- 1. Trayectoria Circular en el Suelo ---
        \draw[dashed, gray, thin] (0,0,0) circle (\L);

        % --- 2. Eje Central (Poste) ---
        % Base y tapa para darle solidez
        \def\Raxis{0.4} % Definimos el radio del poste para que sea consistente

        % Cara inferior (Base)

        % Cuerpo del cilindro (Paredes)
        % Dibujamos un rectángulo que une los extremos laterales del círculo
        \draw[fill=gray!30, draw=black,thick](\Raxis+0.1,-0.1,0) arc (-50:130:\Raxis)--++(0,0,\Hshaft)-- ++(0.87,-0.28,0.21)--cycle;
            
        % Cara superior (Tapa)
        \draw[fill=gray!30, thick] (0,0,\Hshaft-0.15) circle (\Raxis);

        % Línea de eje y punto O
        \draw[dash dot, thin] (0,0,-0.5) -- (0,0,\Hshaft + 1.0);
        \coordinate (O) at (0,0,\zO);

        % --- 3. Brazo Rotatorio ---
        \coordinate (B) at (\L, 0, \zO);
        % Línea gruesa y sólida
        \draw[line width=5pt] (O) -- (B) node[midway, above] {$L$};
        \draw[line width=4pt, gray!70, line cap=round] (O) -- (B);

        \filldraw (O) circle (1.5pt) node[left, xshift=-2pt,  yshift=2pt] {$O$};
        % --- 4. Disco en el punto B (Mejorado) ---
        % Rotamos el sistema local para dibujar el disco vertical
        \tdplotsetrotatedcoords{0}{90}{0}
        \begin{scope}[tdplot_rotated_coords, shift={(B)}]
            % Círculo trasero (para dar profundidad)
            \draw[fill=gray!60, thick] (0,0,-\DiskThick) circle (\R);
            
            % Borde del disco (conecta los círculos)
            \draw[thick, gray!60] (0,\R,0) -- (0,\R,-\DiskThick);
            \draw[thick, gray!60] (0,-\R,0) -- (0,-\R,-\DiskThick);

            % Círculo frontal (cara principal)
            \draw[fill=gray!30, thick] (0,0,0) circle (\R);
            
            % Eje del disco
            \draw[fill=gray!80] (0,0,0) circle (0.15);

            % Puntos y Etiquetas
            \filldraw[black] (0,0,0) circle (1.5pt) node[above right] {$B$};
            % Vector Radio R apuntando al suelo
            \draw[->, thick] (0,0,0) -- (0,-\R,0) node[midway, above] {$R$};
        \end{scope}

        % --- 5. Indicador de Rotación Omega ---
        \draw[->, very thick] (0,0,\Hshaft + 0.4) ++(120:0.6) arc (120:420:0.6);
        \node[left] at (-0.6, 0, \Hshaft + 0.7) {$\Omega$};

        % --- Hack para Centrado ---
        % Puntos invisibles extremos para equilibrar la caja delimitadora
        \path (\L+\R, 0, 0); \path (-\L-\R, 0, 0);

    \end{tikzpicture}
    \caption{Representación mejorada del sistema de rueda pivotante. El disco de radio $R$ rueda sobre una trayectoria circular de radio $L$, impulsado por el brazo que gira con velocidad $\Omega$ alrededor del eje vertical.}
    \label{fig:rueda_brazo_final}
\end{figure}